\documentclass[12pt, a4paper]{article}

\usepackage[a4paper, landscape, top = 1 in, bottom = 1 in, left = 1 in, right = 1 in]{geometry}

\usepackage{amsmath, amssymb, amsfonts}
\allowdisplaybreaks[4]

\usepackage{enumerate}
\usepackage{multirow}
\usepackage{hhline}
\usepackage{array}
\usepackage{longtable}
\usepackage{graphicx}
\usepackage{tabularray}
\usepackage{undertilde}
\usepackage{xcolor}
\usepackage{dingbat}
\usepackage{fontawesome5}
\usepackage[colorlinks=true, linkcolor=blue, urlcolor=red]{hyperref}
\usepackage{tasks}
\usepackage{bbding}
\usepackage{twemojis}
% how to use bull's eye ----- \scalebox{2.0}{\twemoji{bullseye}}
\usepackage{customdice}
% how to put dice face ------ \dice{2}
\usepackage{fontspec}
\usepackage{fancyhdr}
\usepackage{lastpage}



\pagestyle{fancy}
\fancyhf{}
\fancyfoot[C]{Page \thepage\ of \pageref{LastPage}}
\renewcommand{\headrulewidth}{0pt}



\title{MSMS 401 : Bayesian Statistics and Computation}
\author{{\fontsize{25}{20} \benyorithafont Ananda Biswas}}
\date{}

\newfontface\gloriahallelujahfont{GLORIAHALLELUJAH-REGULAR.TTF}

\newfontface\benyorithafont{Benyoritha-G334D.otf}


\begin{document}

\maketitle

\section*{\faArrowAltCircleRight[regular] \textcolor{blue}{Question}}

Consider a random sample of size $n$ from $N(\theta, \sigma_0^2)$, $\sigma_0^2$ is the known variance. Take a conjugate prior for $\theta$. Obtain the posterior distribution of $\theta$, hence obtain Bayes estimator of $\theta$ assuming squared-error loss function.

\section*{\faArrowAltCircleRight[regular] \textcolor{blue}{Answer}}

Let $x_1, x_2, \ldots, x_n \overset{iid}{\sim} f(x \mid \theta)$ where

\begin{align*}
f(x \mid \theta) = \dfrac{1}{\sigma_0 \sqrt{2\pi}} \exp\left\{ -\dfrac{(x-\theta)^2}{2\sigma_0^2} \right\}.
\end{align*}

Consider \textit{a priori} $\theta \sim N(\mu, \tau^2)$ \textit{i.e.}

\begin{align*}
\theta \sim \pi(\theta \mid \mu, \tau^2) = \dfrac{1}{\tau\sqrt{2\pi}} \exp\left\{ -\dfrac{(\theta-\mu)^2}{2\tau^2} \right\}.
\end{align*}

\newpage

The likelihood of $\theta$ is

\begin{align*}
L(\theta)
&= \prod_{i=1}^{n} \dfrac{1}{\sigma_0\sqrt{2\pi}} \exp\!\left\{ -\dfrac{(x_i-\theta)^2}{2\sigma_0^2} \right\} \\[0.8em]
&= \left(\dfrac{1}{2\pi\sigma_0^2}\right)^{n/2} \exp\left\{ -\dfrac{1}{2\sigma_0^2} \sum\limits_{i=1}^{n} (x_i-\theta)^2 \right\}.
\end{align*}

We compute the posterior distribution of $\theta$ as

\begin{align*}
p(\theta \mid \mathbf{x}) = \dfrac{L(\theta) \, \cdot \, \pi(\theta)}{\displaystyle \int\limits_{-\infty}^{\infty} L(\theta) \, \cdot \, \pi(\theta) \, d\theta}.
\end{align*}

\begin{align*}
\therefore\; p(\theta \mid \mathbf{x}) = \dfrac{\left(\dfrac{1}{\sigma_0\sqrt{2\pi}}\right)^n \cdot \dfrac{1}{\tau\sqrt{2\pi}} \cdot \exp\left\{-\dfrac{1}{2\sigma_0^2}\displaystyle\sum\limits_{i=1}^{n}(x_i-\theta)^2\right\} \cdot \exp\left\{-\dfrac{(\theta-\mu)^2}{2\tau^2}\right\}}{\displaystyle \int\limits_{-\infty}^{\infty} \left(\dfrac{1}{\sigma_0\sqrt{2\pi}}\right)^n \cdot \dfrac{1}{\tau\sqrt{2\pi}} \cdot \exp\left\{-\dfrac{1}{2\sigma_0^2}\sum\limits_{i=1}^{n}(x_i-\theta)^2\right\} \cdot \exp\left\{-\dfrac{(\theta-\mu)^2}{2\tau^2}\right\} d\theta }.
\end{align*}

We focus on simplifying the integral in the denominator.

\begin{align*}
\text{Say, } I = {\displaystyle \int\limits_{-\infty}^{\infty} \left(\dfrac{1}{\sigma_0\sqrt{2\pi}}\right)^n \cdot \dfrac{1}{\tau\sqrt{2\pi}} \cdot \exp\left\{-\dfrac{1}{2\sigma_0^2}\sum\limits_{i=1}^{n}(x_i-\theta)^2\right\} \cdot \exp\left\{-\dfrac{(\theta-\mu)^2}{2\tau^2}\right\} d\theta }.
\end{align*}

\newpage

Using the identity $\sum\limits_{i=1}^{n}(x_i-\theta)^2 = \sum\limits_{i=1}^{n}(x_i-\bar{x})^2 + n(\bar{x}-\theta)^2$, we get

\begin{align*}
I &=
\left(\dfrac{1}{\sigma_0\sqrt{2\pi}}\right)^n \cdot \dfrac{1}{\tau\sqrt{2\pi}} \int\limits_{-\infty}^{\infty} \exp\left\{ \textcolor[HTML]{AE28B8}{
-\dfrac{1}{2\sigma_0^2} \sum\limits_{i=1}^{n} (x_i - \bar{x})^2 } \textcolor[HTML]{AE28B8}{- \dfrac{1}{2\sigma_0^2}n(\bar{x}-\theta)^2} - \dfrac{(\theta-\mu)^2}{2\tau^2} \right\} d\theta \\[1em]
&=
\left(\dfrac{1}{\sigma_0\sqrt{2\pi}}\right)^n \cdot \dfrac{1}{\tau\sqrt{2\pi}} \cdot \exp\left\{ \textcolor[HTML]{AE28B8}{-\dfrac{1}{2\sigma_0^2}\sum\limits_{i=1}^{n}(x_i-\bar{x})^2} \right\} \cdot \int\limits_{-\infty}^{\infty} \exp\left\{ \textcolor[HTML]{AE28B8}{- \dfrac{1}{2\sigma_0^2}n(\bar{x}-\theta)^2} - \dfrac{(\theta-\mu)^2}{2\tau^2} \right\} d\theta \\[1em]
&=
\left(\dfrac{1}{\sigma_0\sqrt{2\pi}}\right)^n \cdot \dfrac{1}{\tau\sqrt{2\pi}} \cdot \exp\left\{ -\dfrac{1}{2\sigma_0^2}\sum\limits_{i=1}^{n}(x_i-\bar{x})^2 \right\} \cdot \int\limits_{-\infty}^{\infty} \exp\left\{ -\dfrac{1}{2}\left[ \textcolor[HTML]{D4A420}{\dfrac{n}{\sigma_0^2}(\theta^2-2\bar{x}\theta+\bar{x}^2)} \textcolor[HTML]{53D420}{+ \dfrac{1}{\tau^2}(\theta^2-2\mu\theta+\mu^2)} \right] \right\} d\theta \\[1em]
&=
\left(\dfrac{1}{\sigma_0\sqrt{2\pi}}\right)^n \cdot \dfrac{1}{\tau\sqrt{2\pi}} \cdot \exp\left\{ -\dfrac{1}{2\sigma_0^2}\sum\limits_{i=1}^{n}(x_i-\bar{x})^2 \right\} \cdot \int\limits_{-\infty}^{\infty} \exp\left\{ -\dfrac{1}{2}\left[ \textcolor[HTML]{D45320}{\left(\dfrac{n}{\sigma_0^2}+\dfrac{1}{\tau^2}\right)\theta^2 - 2\left(\dfrac{n\bar{x}}{\sigma_0^2}+\dfrac{\mu}{\tau^2}\right)\theta + \dfrac{n\bar{x}^2}{\sigma_0^2} + \dfrac{\mu^2}{\tau^2}} \right] \right\} d\theta
\end{align*}

To complete the square, we define
\begin{align*}
A &= \dfrac{n}{\sigma_0^2} + \dfrac{1}{\tau^2} \\[0.5em]
B &= \dfrac{n\bar{x}}{\sigma_0^2} + \dfrac{\mu}{\tau^2}.
\end{align*}

$$\therefore A\theta^2 - 2B\theta = A\left(\theta - \dfrac{B}{A}\right)^2 - \dfrac{B^2}{A}.$$

\newpage

\begin{align*}
\therefore I &=
\left(\dfrac{1}{\sigma_0\sqrt{2\pi}}\right)^n \cdot \dfrac{1}{\tau\sqrt{2\pi}} \cdot \exp\left\{ -\dfrac{1}{2\sigma_0^2}\sum\limits_{i=1}^{n}(x_i-\bar{x})^2 \right\} \cdot \int\limits_{-\infty}^{\infty} \exp\left\{ \textcolor[HTML]{205CD4}{-\dfrac{1}{2}A\left(\theta-\dfrac{B}{A}\right)^2 + \dfrac{1}{2}\dfrac{B^2}{A} - \dfrac{1}{2}\left(\dfrac{n\bar{x}^2}{\sigma_0^2} + \dfrac{\mu^2}{\tau^2}\right)} \right\} d\theta \\[1em]
&=
\left(\dfrac{1}{\sigma_0\sqrt{2\pi}}\right)^n \cdot \dfrac{1}{\tau\sqrt{2\pi}} \cdot \exp\left\{ -\dfrac{1}{2\sigma_0^2}\sum\limits_{i=1}^{n}(x_i-\bar{x})^2 \right\} \cdot \exp\left\{ \textcolor[HTML]{205CD4}{\dfrac{1}{2}\dfrac{B^2}{A} - \dfrac{1}{2}\left(\dfrac{n\bar{x}^2}{\sigma_0^2} + \dfrac{\mu^2}{\tau^2}\right)} \right\} \cdot \int\limits_{-\infty}^{\infty} \exp\left\{ \textcolor[HTML]{205CD4}{-\dfrac{1}{2}A\left(\theta-\dfrac{B}{A}\right)^2} \right\} d\theta \\[1em]
&=
\left(\dfrac{1}{\sigma_0\sqrt{2\pi}}\right)^n \cdot \dfrac{1}{\tau\sqrt{2\pi}} \cdot \exp\left\{ -\dfrac{1}{2\sigma_0^2}\sum\limits_{i=1}^{n}(x_i-\bar{x})^2 \right\} \cdot \exp\left\{ \dfrac{1}{2}\dfrac{B^2}{A} - \dfrac{1}{2}\left(\dfrac{n\bar{x}^2}{\sigma_0^2} + \dfrac{\mu^2}{\tau^2}\right) \right\} \cdot \textcolor[HTML]{C71056}{\sqrt{\dfrac{2\pi}{A}}}
\end{align*}

Notice that the numerator is same, only it does not have the integral. So a similar simplification of the numerator will result in

\begin{align*}
\left(\dfrac{1}{\sigma_0\sqrt{2\pi}}\right)^n \cdot \dfrac{1}{\tau\sqrt{2\pi}} \cdot \exp\left\{ -\dfrac{1}{2\sigma_0^2}\sum\limits_{i=1}^{n}(x_i-\bar{x})^2 \right\} \cdot \exp\left\{ \dfrac{1}{2}\dfrac{B^2}{A} - \dfrac{1}{2}\left(\dfrac{n\bar{x}^2}{\sigma_0^2} + \dfrac{\mu^2}{\tau^2}\right) \right\} \cdot \exp\left\{ -\dfrac{1}{2}A\left(\theta-\dfrac{B}{A}\right)^2 \right\}.
\end{align*}

Then,

\begin{align*}
p(\theta \mid \mathbf{x}) &= \dfrac{\textcolor[HTML]{EB0505}{\left(\dfrac{1}{\sigma_0\sqrt{2\pi}}\right)^n \cdot \dfrac{1}{\tau\sqrt{2\pi}} \cdot \exp\left\{ -\dfrac{1}{2\sigma_0^2}\sum\limits_{i=1}^{n}(x_i-\bar{x})^2 \right\} \cdot \exp\left\{ \dfrac{1}{2}\dfrac{B^2}{A} - \dfrac{1}{2}\left(\dfrac{n\bar{x}^2}{\sigma_0^2} + \dfrac{\mu^2}{\tau^2}\right) \right\}} \cdot \exp\left\{ -\dfrac{1}{2}A\left(\theta-\dfrac{B}{A}\right)^2 \right\}}{\textcolor[HTML]{EB0505}{\left(\dfrac{1}{\sigma_0\sqrt{2\pi}}\right)^n \cdot \dfrac{1}{\tau\sqrt{2\pi}} \cdot \exp\left\{ -\dfrac{1}{2\sigma_0^2}\sum\limits_{i=1}^{n}(x_i-\bar{x})^2 \right\} \cdot \exp\left\{ \dfrac{1}{2}\dfrac{B^2}{A} - \dfrac{1}{2}\left(\dfrac{n\bar{x}^2}{\sigma_0^2} + \dfrac{\mu^2}{\tau^2}\right) \right\}} \cdot \sqrt{\dfrac{2\pi}{A}}} \\[1em]
&=
\sqrt{\dfrac{A}{2\pi}} \cdot \exp\left\{ -\dfrac{1}{2}A\left(\theta-\dfrac{B}{A}\right)^2 \right\}
\end{align*}

which we right away identify as the PDF of $N\left( \dfrac{B}{A}, \dfrac{1}{A}\right)$.

\newpage

Thus, \textit{a posteriori}

\begin{align*}
\theta \mid \mathbf{x} \sim N\left( \dfrac{\dfrac{n\bar{x}}{\sigma_0^2} + \dfrac{\mu}{\tau^2}}{\dfrac{n}{\sigma_0^2} + \dfrac{1}{\tau^2}}, \dfrac{1}{\dfrac{n}{\sigma_0^2} + \dfrac{1}{\tau^2}} \right).
\end{align*}
\end{document}